%============================================================
% PREÁMBULO OFICIAL FLUIDODINAMICA_BASE
%============================================================

\documentclass[12pt, a4paper]{article}

%------------------------------------------------------------
% Idioma, codificación y tipografía
%------------------------------------------------------------
\usepackage[spanish]{babel}
\usepackage[utf8]{inputenc}
\usepackage[T1]{fontenc}
\usepackage{lmodern}        % Tipografía moderna, clara y apta para PDF
\usepackage{microtype}      % Optimiza microtipografía (ideal para libros y apuntes)

%------------------------------------------------------------
% Matemática
%------------------------------------------------------------
\usepackage{amsmath, amssymb, amsfonts}
\usepackage{bm}             % Vectores y tensores en negrita: \bm{u}
\usepackage{mathtools}      % Extensiones de amsmath
\usepackage{siunitx}        % Notación correcta de unidades SI
\sisetup{
  locale = DE,            % usa coma decimal si querés; cambiar a US para punto
  per-mode = symbol,
  exponent-product = \cdot,
  separate-uncertainty = true
}

%------------------------------------------------------------
% Gráficos y figuras
%------------------------------------------------------------
\usepackage{graphicx}
\usepackage{float}
\usepackage{subcaption}     % Figuras con subfiguras
\usepackage{wrapfig}

%------------------------------------------------------------
% Diseño de página
%------------------------------------------------------------
\usepackage{geometry}
\geometry{
  a4paper,
  left=25mm,
  right=25mm,
  top=25mm,
  bottom=25mm
}

\usepackage{setspace}       % Control de espaciado (opcional)
\onehalfspacing             % Espaciado 1.5 (queda excelente para apuntes)

%------------------------------------------------------------
% Tablas profesionales
%------------------------------------------------------------
\usepackage{booktabs}       % Líneas horizontales profesionales
\usepackage{multirow}
\usepackage{tabularx}
\usepackage{longtable}

%------------------------------------------------------------
% Enlaces
%------------------------------------------------------------
\usepackage[hidelinks]{hyperref}  % Enlaces sin recuadros

%------------------------------------------------------------
% Colores y cajas
%------------------------------------------------------------
\usepackage{xcolor}
\usepackage{tcolorbox}
\tcbset{
  colback = white,
  colframe = black,
  arc = 2mm,
  boxrule = 0.6pt
}

%------------------------------------------------------------
% Ambientes personalizados
%------------------------------------------------------------

% Entorno para ejercicios
\newenvironment{ejercicio}[1][]{
  \begin{tcolorbox}[title={\textbf{Ejercicio #1}}, colback=gray!5]
}{
  \end{tcolorbox}
}

% Entorno para soluciones
\newenvironment{solucion}{
  \begin{tcolorbox}[title={\textbf{Solución}}, colback=blue!3]
}{
  \end{tcolorbox}
}

% Para dejar comentarios o notas al margen
\newcommand{\nota}[1]{\textcolor{blue}{\textbf{#1}}}

%============================================================
% FIN DEL PREÁMBULO
%============================================================
\begin{document}

%============================================================
% Ejercicio p.2 – Tubo Pitot en cañería de 1 1/2"
%============================================================

\subsection*{Ejercicio p.3: Medición de caudal con tubo Pitot}

Por una cañería de acero comercial IPS Sch 40 de diámetro nominal
$1{\,}1/2''$ en la que circula agua se encuentra instalado un tubo Pitot.
Si el tubo queda ubicado en el centro de la cañería, determinar el flujo
másico de agua (en kg/s) cuando la lectura del manómetro asociado al
tubo Pitot es de $6{,}1\,\text{kPa}$.

\medskip

\noindent
\textbf{Datos adicionales:}
\begin{itemize}
  \item Agua: $\rho = 1000\,\text{kg/m}^3$, $\mu = 1\times 10^{-3}\,\text{Pa\,s}$.
  \item Cañería IPS Sch 40, $\varnothing_N = 1{\,}1/2''$:
        diámetro interno $D \simeq 0{,}0408\,\text{m}$.
  \item Lectura del manómetro del Pitot:
        $\Delta P = 6{,}1\,\text{kPa} = 6{,}1\times 10^3\,\text{Pa}$.
\end{itemize}

\noindent
\textbf{Incógnita:} flujo másico de agua $\dot m$ [kg/s].

%------------------------------------------------------------
\subsubsection*{1. Hipótesis}

\begin{enumerate}
  \item Flujo estacionario de agua, incompresible.
  \item Cañería de sección circular constante; flujo plenamente desarrollado.
  \item Tubo Pitot simple con coeficiente de recuperación de presión
        $C_p \approx 1$ (como en las notas de clase de medidores).
  \item El Pitot está ubicado en el centro de la cañería, de modo que mide
        la \emph{velocidad local máxima} $U_{\max}$.
  \item Se supone que el flujo es completamente turbulento en la cañería,
        por lo que el cociente entre velocidad media y máxima
        $\langle U\rangle/U_{\max}$ se toma de la curva de la página
        correspondiente del teórico (relación de velocidades en función
        del número de Reynolds).
\end{enumerate}

%------------------------------------------------------------
\subsubsection*{2. Modelo del tubo Pitot}

\paragraph{2.1. Relación entre $\Delta P$ y la velocidad máxima.}

El tubo Pitot mide la diferencia entre presión total y estática en el punto
donde está inmerso. Despreciando pérdidas internas en el Pitot e introduciendo
un coeficiente $C_p \approx 1$, la relación entre la lectura de presión
y la velocidad local máxima $U_{\max}$ es

\[
\Delta P = \frac{1}{2}\,\rho\,U_{\max}^2
\quad\Longrightarrow\quad
U_{\max} = \sqrt{\frac{2\,\Delta P}{\rho}}.
\]

Sustituyendo:

\[
U_{\max}
= \sqrt{\frac{2\cdot 6{,}1\times 10^3}{1000}}
= \sqrt{12{,}2}
\approx 3{,}49\,\text{m/s}.
\]

\paragraph{2.2. Cálculo del número de Reynolds.}

Primero evaluamos el número de Reynolds basado en la velocidad máxima:

\[
\mathrm{Re}_{\max}
= \frac{\rho\,U_{\max} D}{\mu}
= \frac{1000\cdot 3{,}49\cdot 0{,}0408}{1\times 10^{-3}}
\approx 1{,}4\times 10^5.
\]

Este valor se sitúa en la región completamente turbulenta de la curva
$\langle U\rangle/U_{\max}$ vs. $(D\,U_{\max}\rho/\mu)$ incluida en las
notas de clase (medidores\_caudal, págs.~29--35).

\paragraph{2.3. Relación entre velocidad media y máxima.}

De la gráfica ``Relación de velocidades en función del número de Reynolds''
para $D\,U_{\max}\rho/\mu \sim 10^5$ se lee:

\[
\frac{\langle U\rangle}{U_{\max}} \simeq 0{,}83.
\]

Por tanto, la velocidad media en la cañería es

\[
\langle U\rangle = 0{,}83\,U_{\max}
\approx 0{,}83\cdot 3{,}49
\approx 2{,}90\,\text{m/s}.
\]

%------------------------------------------------------------
\subsubsection*{3. Caudal volumétrico y flujo másico}

El área de la sección de la cañería es

\[
A = \frac{\pi D^2}{4}
  = \frac{\pi (0{,}0408)^2}{4}
  \approx 1{,}31\times 10^{-3}\,\text{m}^2.
\]

El caudal volumétrico resulta

\[
Q = \langle U\rangle\,A
  \approx 2{,}90\cdot 1{,}31\times 10^{-3}
  \approx 3{,}8\times 10^{-3}\,\text{m}^3/\text{s}.
\]

El flujo másico de agua es

\[
\dot m = \rho\,Q
       = 1000\cdot 3{,}8\times 10^{-3}
       \approx 3{,}8\,\text{kg/s}.
\]

\[
\boxed{\dot m \approx 3{,}8\ \text{kg/s}.}
\]

%------------------------------------------------------------
\subsubsection*{4. Verificación de coherencia física}

\paragraph{Número de Reynolds con velocidad media.}

Con la velocidad media obtenida:

\[
\mathrm{Re} = \frac{\rho\,\langle U\rangle D}{\mu}
            = \frac{1000\cdot 2{,}90\cdot 0{,}0408}{1\times 10^{-3}}
            \approx 1{,}2\times 10^5,
\]

que confirma el régimen plenamente turbulento, consistente con el uso
de la relación $\langle U\rangle/U_{\max}\simeq 0{,}83$ tomada de la curva.

\paragraph{Comentario final.}

El procedimiento seguido es exactamente el del teórico de medidores:
\begin{enumerate}
  \item se obtiene $U_{\max}$ a partir de la lectura del Pitot,
  \item se calcula $\mathrm{Re}_{\max}$ y se localiza el punto en la gráfica
        $\langle U\rangle/U_{\max}$ vs. $D\,U_{\max}\rho/\mu$,
  \item se determina la velocidad media a partir de 
        $\langle U\rangle/U_{\max}$,
  \item y finalmente se halla el flujo másico 
        $\dot m = \rho\,\langle U\rangle A$.
\end{enumerate}
El valor $\dot m \approx 3{,}8\,\text{kg/s}$ es razonable para una cañería
de $1{\,}1/2''$ con una diferencia de presión de Pitot del orden de varios kPa.

\end{document}